\DocumentMetadata{
  pdfversion=2.0,
  pdfstandard=ua-2,
  lang=en-US,
  tagging=on,
  tagging-setup={math/setup=mathml-SE,tabsorder=structure}
}
\documentclass[11pt,letterpaper]{article}
\usepackage[margin=0.68in,includefoot]{geometry}
\usepackage{newcomputermodern}
\usepackage{amsmath,amscd,mathtools}
\usepackage{tikz,adjustbox}
\providecommand{\square}{\mdlgwhtsquare}
\usepackage[hidelinks]{hyperref}
\hypersetup{
  pdftitle={Topology First-Year Exam, May 2013, Part 1},
  pdfauthor={University of Florida Department of Mathematics},
  pdfsubject={Graduate mathematics examination},
  pdfdisplaydoctitle=true,
  bookmarksopen=true
}
\setcounter{secnumdepth}{0}
\setlength{\parindent}{0pt}
\setlength{\parskip}{0.28em}
\setlength{\emergencystretch}{3em}
\setlength{\leftmargini}{2.15em}
\setlength{\leftmarginii}{2.65em}
\setlength{\leftmarginiii}{3em}
\pagestyle{empty}
\raggedbottom
\allowdisplaybreaks
\NewTaggingSocketPlug{block/list/label}{examproblem}{%
  \tagstructbegin{tag=Lbl}\tagstructend%
  \tagstructbegin{tag=\LBody}%
  \tagstructbegin{tag=\ProblemHeadingTag,title-o={\ProblemHeadingText}}%
  \tagmcbegin{tag=\ProblemHeadingTag,actualtext={\ProblemHeadingText}}%
  #2%
  \tagmcend%
  \tagstructend%
}
\newenvironment{examproblems}[2]{%
  \def\ProblemHeadingTag{#2}%
  \AssignTaggingSocketPlug{block/list/label}{examproblem}%
  \begin{enumerate}%
  \setcounter{enumi}{#1}%
  \renewcommand{\labelenumi}{\textbf{\theenumi.}}%
  \setlength{\topsep}{0.35em}%
  \setlength{\partopsep}{0pt}%
  \setlength{\itemsep}{0.48em}%
  \setlength{\parsep}{0.18em}%
}{\end{enumerate}}
\newenvironment{examsubparts}{%
  \AssignTaggingSocketPlug{block/list/label}{default}%
  \begin{enumerate}%
  \setlength{\topsep}{0.18em}%
  \setlength{\partopsep}{0pt}%
  \setlength{\itemsep}{0.16em}%
  \setlength{\parsep}{0.1em}%
}{\end{enumerate}}
\newenvironment{examinstructions}{%
  \par\begingroup\itshape
}{%
  \par\endgroup\vspace{0.18em}
}
\makeatletter
\renewcommand\subsection{\@startsection{subsection}{2}{0pt}%
  {-1.25ex plus -.25ex minus -.1ex}%
  {0.55ex plus .1ex}%
  {\normalfont\large\bfseries}}
\renewcommand{\@maketitle}{%
  \newpage\null\vskip -1em%
  {\footnotesize\bfseries
    \href{https://www.ufl.edu/}{University of Florida}, \href{https://math.ufl.edu/}{Department of Mathematics}\par}%
  \vskip -0.5em%
  \tagstructbegin{tag=H1,title={Topology First-Year Exam, May 2013, Part 1}}%
  \tagpdfparaOff%
  \tagmcbegin{tag=H1}%
    {\LARGE\bfseries\href{https://gma.math.ufl.edu/exams/topology-first-year/}{Topology First-Year Exam}, May 2013, Part 1\par}%
  \tagmcend%
  \tagpdfparaOn%
  \tagstructend%
  \vskip 0.25em%
  \tagmcbegin{artifact}%
    \hrule height 1pt%
  \tagmcend%
  \vskip 1em%
}
\makeatother
\title{Topology First-Year Exam, May 2013, Part 1}
\author{}
\date{}
\begin{document}
\maketitle
\thispagestyle{empty}
\begin{examinstructions}
Work the following problems and show all work. Support all statements to the best of your ability.
\end{examinstructions}
\begin{examproblems}{0}{H2}
\def\ProblemHeadingText{Problem 1.}
\item
Prove that the spaces \([0,1]\) and \([0,1)\) are not homeomorphic.
\def\ProblemHeadingText{Problem 2.}
\item
Let~\(X\) be a connected normal topological space having more than one point. Prove that~\(X\) is uncountable.
\def\ProblemHeadingText{Problem 3.}
\item
Let \(f:X\to Y\) be a continuous map to a Hausdorff space. Prove that the graph of~\(f\), \(G=\{(x,f(x))\mid x\in X\}\), is a closed subset of \(X\times Y\).
\def\ProblemHeadingText{Problem 4.}
\item
This problem has two parts.
\begin{examsubparts}
\item[\textnormal{(a)}]
Let~\(X\) be a compact space. Show that for every topological imbedding \(f:X\to Y\) into a Hausdorff space the image \(f(X)\) is closed in~\(Y\).
\item[\textnormal{(b)}]
Suppose that a normal topological space~\(X\) has the property that for every topological imbedding \(f:X\to Y\), the image \(f(X)\) is closed in~\(Y\). Does it follow that~\(X\) is compact?
\end{examsubparts}
\def\ProblemHeadingText{Problem 5.}
\item
Is every closed subset~\(A\) of a separable space~\(X\) separable itself if
\begin{examsubparts}
\item[\textnormal{(a)}]
\(X=\mathbb{R}\)?
\item[\textnormal{(b)}]
\(X=\mathbb{R}\times\mathbb{R}\)?
\item[\textnormal{(c)}]
\(X=\mathbb{R}_\ell\)?
\item[\textnormal{(d)}]
\(X=\mathbb{R}_\ell\times\mathbb{R}_\ell\)?
\end{examsubparts}
\def\ProblemHeadingText{Problem 6.}
\item
Does there exist a covering map \(p:\mathbb{R}^2\to\mathbb{RP}^2\) from the Euclidean plane to the projective plane?
\end{examproblems}
\begin{examinstructions}
Answer the following with complete definitions or statements or short proofs.
\end{examinstructions}
\begin{examproblems}{6}{H2}
\def\ProblemHeadingText{Problem 7.}
\item
State the Tietze Extension Theorem.
\def\ProblemHeadingText{Problem 8.}
\item
Is the space \(\mathbb{R}^\omega\) connected in the uniform topology?
\def\ProblemHeadingText{Problem 9.}
\item
Does there exist a continuous surjective map from the 2-sphere \(S^2\) to the punctured square \(([-1,1]\times[-1,1])\setminus\{(0,0)\}\)?
\def\ProblemHeadingText{Problem 10.}
\item
Is every connected space path connected?
\def\ProblemHeadingText{Problem 11.}
\item
What is a basis of a topology? What is a subbasis?
\def\ProblemHeadingText{Problem 12.}
\item
State the Baire Category Theorem.
\def\ProblemHeadingText{Problem 13.}
\item
Is the unit circle \(S^1=\{x\in\mathbb{R}^2\mid \lVert x\rVert=1\}\) a
\begin{examsubparts}
\item[\textnormal{(a)}]
retract of \(\mathbb{R}^2\)?
\item[\textnormal{(b)}]
retract of \(\mathbb{R}^2\setminus\{(2,0)\}\)?
\item[\textnormal{(c)}]
deformation retract of \(\mathbb{R}^2\setminus\{(0,0)\}\)?
\item[\textnormal{(d)}]
deformation retract of \(\mathbb{R}^2\setminus\{(0,0),(2,0)\}\)?
\item[\textnormal{(e)}]
retract of \(\mathbb{R}^2\setminus\{(0,0),(2,0)\}\)?
\end{examsubparts}
\def\ProblemHeadingText{Problem 14.}
\item
State the Brouwer fixed point theorem. Does every continuous map \(f:[0,1]\times[0,1]\to[0,1)\times[0,1)\) have a fixed point?
\def\ProblemHeadingText{Problem 15.}
\item
Can the space of irrationals in the subspace topology be presented as a countable union of nowhere dense subsets?
\end{examproblems}
\end{document}
